\chapter{Parity, Benchmarks, and Protocols}

\ChapterMeta{
This chapter details parity checking methodologies, protocol-pinned evaluation configurations, and the utilization of benchmark and sweep harnesses.
}{
These tools ensure that comparisons remain stable and actionable by proactively detecting silent numerical drift and systematically tracking performance and metric trends over time.
}{
Execution requires a minimum of two comparable prediction artifacts for parity analysis; protocol-driven runs necessitate pinned settings files alongside consistent preprocessing and output formats.
}

\section{Backend parity}
When exporting predictions across disparate backends (e.g., PyTorch, ONNX Runtime, TensorRT), it is critical to detect any silent numerical drift. Parity tooling systematically compares two prediction artifacts and generates a comprehensive report of any discrepancies.

\begin{figure}[h]
\centering
\resizebox{\linewidth}{!}{%
\begin{tikzpicture}[
  node distance=10mm,
  font=\small,
  box/.style={draw, rounded corners, align=left, inner sep=6pt},
  arrow/.style={-{Latex[length=2mm]}, thick},
  dashedarrow/.style={-{Latex[length=2mm]}, thick, dashed}
]

\node[box] (inputs) {\textbf{Fixed inputs}\\same image subset\\same preprocessing};

\node[box, right=18mm of inputs, yshift=10mm] (a) {\textbf{Backend A}\\torch / onnxrt};
\node[box, right=18mm of inputs, yshift=-10mm] (b) {\textbf{Backend B}\\trt / opencv-dnn};

\node[box, right=20mm of a] (pa) {\path{predictions\_a.json}};
\node[box, right=20mm of b] (pb) {\path{predictions\_b.json}};

\node[box, below=12mm of pb] (par) {\textbf{Parity report}\\diff stats\\(\path{parity.json})};
\node[box, right=20mm of par] (eval) {\textbf{Eval reports}\\metrics JSON\\protocol-pinned};

\draw[dashedarrow] (inputs) -- (a);
\draw[dashedarrow] (inputs) -- (b);
\draw[arrow] (a) -- (pa);
\draw[arrow] (b) -- (pb);
\draw[arrow] (pa) -- (par);
\draw[arrow] (pb) -- (par);
\draw[arrow] (pa) -- (eval);
\draw[arrow] (pb) -- (eval);

\end{tikzpicture}
}
\caption{Apples-to-apples parity: keep inputs fixed, compare exported artifacts, and score with the same evaluator/protocol.}
\end{figure}

\textbf{Conceptual (not copy-paste as-is):} parity command shape.
For a runnable shortest command, use Chapter~5 Workflow A with real artifact paths.
\begin{lstlisting}[language=bash]
yolozu parity   --reference reports/pred_torch.json   --candidate reports/pred_onnxrt.json
\end{lstlisting}

\section{Protocols (pinning evaluation settings)}
Protocols establish immutable evaluation configurations (such as dataset splits, image dimensions, and specific metric keys). This standardization is paramount for ensuring stable, longitudinal comparisons.

For YOLO26 comparisons, the pinned protocol is defined as follows:
\begin{itemize}
  \item Protocol file: \path{protocols/yolo26_eval.json}
  \item Task: \texttt{detect}, Dataset: \texttt{COCO}, Split: \texttt{val2017}
  \item Image size: \texttt{640} (enforcement is the responsibility of the export pipeline)
  \item Bounding box format: \texttt{cxcywh\_norm}
  \item Primary score key: \texttt{map50\_95}
\end{itemize}

A critical semantic rule derived from the protocol documentation:
\begin{itemize}
  \item "End-to-end (e2e) mAP" dictates that the evaluator scores detections exactly as provided; it explicitly does \textit{not} apply Non-Maximum Suppression (NMS).
\end{itemize}

\section{Benchmarks}
Latency and FPS benchmarks constitute a vital component of the regression testing surface.

The recommended user-facing entrypoint is now \cmd{yolozu benchmark}. The
current implementation keeps the honest artifact-first behavior from Phase~1,
and additionally wires real backend orchestration for \cmd{torch},
\cmd{onnx}, and \cmd{engine}. It also accepts \cmd{torchscript} as a
first-class benchmark format, while still recording honest synthetic/skipped
semantics until a dedicated real-orchestration path lands. When the real
backends can execute, the command records a real dataset-pass wall-clock
measurement. When they cannot, it still records skipped formats honestly rather
than fabricating success.

A representative repository-wrapper invocation is:
\begin{lstlisting}[language=bash]
python3 tools/benchmark_model.py \
  --model runs/example/model.pt \
  --onnx-model exports/example.onnx \
  --engine-model exports/example.plan \
  --data data/coco8.yaml \
  --format torch,onnx,engine \
  --protocol nms_applied \
  --latency-source auto \
  --output reports/benchmark_report.json
\end{lstlisting}

The primary artifact set is:
\begin{itemize}
  \item \path{reports/benchmark_report.json}
  \item \path{reports/export_settings_<format>.json}
  \item \path{reports/predictions_<format>.json}
  \item \path{reports/eval_<format>.json}
  \item \path{reports/parity_<format>.json}
\end{itemize}

When \cmd{torch} plus at least one additional real backend (\cmd{onnx} or
\cmd{engine}) complete successfully, \path{parity_<format>.json} is now a real
backend-diff artifact rather than a placeholder. The chosen reference backend
is recorded explicitly (preferring \cmd{torch} when available). Placeholder
parity artifacts remain only for dry-run, skipped, or synthetic-only formats.

\subsection{Gap against the Ultralytics docs surface}
Relative to the public Ultralytics benchmark/export documentation, YOLOZU has
already matched the core benchmark entrypoint shape, but it still trails in
surface breadth.

The most important remaining gaps are:
\begin{itemize}
  \item broader format coverage such as \cmd{openvino}, \cmd{coreml},
  \cmd{saved\_model}, \cmd{tflite}, \cmd{ncnn}, \cmd{rknn}, and \cmd{paddle}
  \item clearer task-specific benchmark expectations for
  \cmd{segmentation}, \cmd{classification}, and \cmd{obb}
  \item stronger validation of format-specific knobs so unsupported
  combinations fail early rather than acting like inert flags
\end{itemize}

This comparison is intentionally about interface behavior and user-facing
surface only. YOLOZU keeps its Apache-2.0-only repository operations policy and
must not pull GPL-licensed implementation code into the repository while
closing these gaps.

For repository-side orchestration, the equivalent wrapper is
\path{tools/benchmark_model.py}.

The benchmark harness, located at \path{tools/benchmark_latency.py}, generates:
\begin{itemize}
  \item A stable JSON report (detailing a single run).
  \item A JSONL history file (facilitating append-only trend tracking).
\end{itemize}

This harness supports both synthetic timing evaluations and configuration-driven, per-bucket executions. In the latter mode, each bucket can ingest externally measured metrics via JSON (for instance, data generated by TensorRT latency measurement scripts).

Typically tracked metrics include FPS and latency quantiles (mean, p50, p90, p95, and p99).

\section{Sweeps}
A sweep harness is indispensable for systematically comparing numerous runs or configurations.

The sweep runner, \path{tools/hpo_sweep.py}, is entirely framework-agnostic and operates by executing templated commands. The core configuration model comprises:
\begin{itemize}
  \item \texttt{base\_cmd}: The command template containing \texttt{\{param\}} placeholders.
  \item \texttt{param\_grid} or \texttt{param\_list}: The defined hyperparameter search space.
  \item Optional keys for metrics extraction and designated output destinations.
\end{itemize}

The default artifacts are designed for effortless diffing and publication:
\path{reports/hpo_sweep.jsonl}, \path{reports/hpo_sweep.csv}, and \path{reports/hpo_sweep.md}.
Utilize the \cmd{--resume} flag to intelligently bypass previously completed parameter combinations.
